\section{Analytic foundations of the two Goldbach arguments}

The two arguments need more than a supply of primes.
They need primes distributed in residue classes, with errors that remain small
when summed against the actual sieve coefficients.
They also need a normalization of the resulting Euler products that is uniform
in the even integer being represented.
These are separate tasks: neither a linear sieve nor an equidistribution
estimate alone gives the final representation count.

\subsection{The main scale and its local densities}

Throughout this section, $N$ is a positive even integer tending to infinity,
$p$ denotes a prime, $\varphi$ is Euler's totient, and $\omega(d)$ counts the
\emph{distinct} prime divisors of $d$.
Write $\gamma$ for the Euler--Mascheroni constant and put
\[
 C_2=\prod_{p>2}\left(1-\frac1{(p-1)^2}\right),\qquad
 \mathfrak S_{\mathrm{Liu}}(N)=C_2\prod_{\substack{p\mid N\\p>2}}\frac{p-1}{p-2},
 \qquad \mathcal X_N=\mathfrak S_{\mathrm{Liu}}(N)\frac{N}{(\log N)^2}.
\]
Here $\mathfrak S_{\mathrm{Liu}}$ is Liu's normalization: the factor at $2$ is omitted.
The source proves $C_2>0$ and $\mathfrak S_{\mathrm{Liu}}(N)\ge C_2$.
Consequently a proved error $O(N/(\log N)^{2+u})$, with fixed $u>0$,
can eventually be made smaller than any prescribed positive multiple of $\mathcal X_N$.
This last absorption is elementary; producing that error is the analytic work.

For the Goldbach difference sequence $N-p$, divisibility by a prime $q\nmid N$
asks for $p\equiv N\pmod q$ among the reduced residue classes.
Its density is therefore $g(q)=1/(q-1)$, extended multiplicatively on squarefree
sifting divisors, not $1/q$.
For an integer cutoff $Z$, the actual sieve product is
\[
 V_N(Z)=\prod_{\substack{p<Z\\p\nmid N}}\left(1-\frac1{p-1}\right).
\]
Evenness removes the exceptional prime $2$.
The literal Jurkat--Richert $1/p$ specialization is a different source model.
What passes between the models is a constructed system of sieve functions and
proved comparison estimates, not an identification of their Euler factors.

\subsection{Prime counting, Mertens estimates, and moving products}

The von Mangoldt function is $\Lambda(p^k)=\log p$ for a prime $p$
and an integer $k\ge1$, and is zero on all other positive integers.
Define $\psi(x)=\sum_{1\le n\le x}\Lambda(n)$,
$\vartheta(x)=\sum_{p\le x}\log p$, and
$\pi(x)=\#\{p\le x:p\text{ is prime}\}$.
The prime-distribution facade supplies a quantitative prime number theorem (PNT):
\inputleannode{bp:foundation-pnt}
For some $c>0$, its error is $O(x\exp(-c(\log x)^{1/10}))$.
The effective $\vartheta$ estimate keeps the prime-power correction explicitly;
Abel summation then supplies prime reciprocal sums and the Mertens estimates.
The separate prime-counting facade also supplies $\pi(n)\sim n/\log n$ and
an upper bound valid for every integer $n\ge2$.

The constant in the product formula is identified, not left unspecified:
\inputleannode{bp:foundation-mertens}
\[
 \left|\prod_{p\le n}(1-1/p)-\frac{e^{-\gamma}}{\log n}\right|
 \le \frac{C}{(\log n)^2}\qquad(n\ge2).
\]
The local Goldbach version also proves a dimension-one interval estimate.
For any finite set $\mathcal P$ of odd primes in $[z_1,z_2)$, with
$2\le z_1\le z_2$, one constant $K>1$ gives
\inputleannode{bp:foundation-local-product}
\[
 \prod_{p\in\mathcal P}\left(1-\frac1{p-1}\right)^{-1}
 \le \frac{\log z_2}{\log z_1}\left(1+\frac K{\log z_1}\right).
\]
This supplies the local-product hypothesis of the actual Chen and Li--Liu
bounding sieves, uniformly over their changing finite prime sets.
Here and below, Li--Liu's $S_1$ means the base sift of the original difference
source $N-p$, and $S_2$ means the sum of sources obtained by fixing one prime
factor; their precise fibres are defined in Section~\ref{sec:liliu-fibres}.
The source $\mathcal B_{10}$ retains three prime labels while sifting the
output $N-rsq$; see Section~\ref{sec:liliu-output-sift}.
These forward references specify the counting objects; the present chapter
explains the analytic estimates that will be applied to them.

Finite truncations of $\mathfrak S_{\mathrm{Liu}}(N)$ must still be distinguished from the full
product. The universal odd-prime tail is independent of $N$; omitted divisor
corrections are not. The source obtains a uniform upper comparison of the
truncation with $(1+\eta)\mathfrak S_{\mathrm{Liu}}(N)$ as the cutoff grows.
The lower comparison needed for $S_1$ requires additional geometry:
for $z\ge N^{1/18}$, there are at most eighteen prime divisors of $N$
at or above $\lceil z\rceil$.
Their correction product is bounded by
$\exp(18/(\lceil z\rceil-2))$ once the denominator is positive.
This is why the lower normalization is not just the upper argument reversed.
The resulting upper and lower products have leading term
$2e^{-\gamma}\mathfrak S_{\mathrm{Liu}}(N)/\log z$, with explicitly paid relative errors.

\subsection{Ordinary prime distribution and the weighted Chen error}

For a real additive normalization $\kappa$ and $x\ge0$, use
\[
 L_\kappa(x)=
 \begin{cases}
 \kappa,&0\le x\le1,\\
 \displaystyle\kappa+\int_2^x\frac{dt}{\log t},&x>1,
 \end{cases}
 \qquad L_*=L_{2/\log2}.
\]
The integral in the second branch is an ordinary oriented integral on an
interval avoiding $1$. The first branch records the implementation's
\emph{totalized interval-integral convention}: an integral over the
nonintegrable logarithmic singularity is assigned zero. It is not a Cauchy
principal value. In particular $L_*(0)=L_*(1)=2/\log2$.
The source's \texttt{trueLogarithmicIntegral} means this $L_*$, not the
unshifted integral and not $x/\log x$; see
\sourcefile{MathlibNt/SieveTheory/Arithmetic/LiuLogarithmicIntegral.lean}{logarithmic-integral normalization}
and \sourcefile{MathlibNt/SieveTheory/Distribution/LiuPan/LiuPanCombinedAbelDeterministic.lean}{the zero-integral endpoint lemma}.
Let $\pi(y;q,a)$ count primes at most the integer $y$ in the class $a\bmod q$,
and, for integers $N\ge0$ and $q\ge1$, define the nested maximum
\[
 E^*(N,q)=\max_{\substack{y\in\mathbb Z\\0\le y\le N}}\ \max_{\substack{0\le a<q\\(a,q)=1}}
 \left|\pi(y;q,a)-\frac{L_*(y)}{\varphi(q)}\right|,
 \qquad Q_B(N)=\left\lfloor\frac{N^{1/2}}{(\log N)^B}\right\rfloor.
\]
The residues $a$ are integers. At modulus zero put $E^*(N,0)=0$;
modulus one has the single reduced residue $a=0$.
Thus the maximum really includes $y=0,1$, with the convention just specified;
see \sourcefile{MathlibNt/SieveTheory/Distribution/BombieriVinogradov.lean}{the exact AP and prefix maxima}.
The proved standard Bombieri--Vinogradov (BV) theorem states that for every $A>0$
there are $B\ge0$ and $C>0$ such that, eventually,
\inputleannode{bp:foundation-bv}
\[
 \sum_{1\le q\le Q_B(N)} E^*(N,q)\le C\frac{N}{(\log N)^A}.
\]
The maxima allow the consumer to choose its residue and its integer endpoint;
the theorem does not bound an arbitrary convolution by ordinary prime BV.
Richert's unshifted-integral formulation is supplied by a further normalization
comparison in the same analytic producer, not by redefining $L_*$.

For comparison with the Chen chapter, its notations $U$ and
$\operatorname{li}_\kappa$ mean $C_2$ and $L_\kappa$, respectively.
Its fixed-endpoint prime error is
\[
 E_{\rm prime}(N,q)=\max_{\substack{0\le a<q\\(a,q)=1}}
 \left|\pi(N;q,a)-\frac{L_*(N)}{\varphi(q)}\right|
 \le E^*(N,q) \qquad(q\ge1).
\]
This is the same centering and residue maximum, with only the additional
prefix maximum omitted.

\paragraph{Why Rosser weights give a lower or upper sieve.}
Let $\mathcal P$ be a finite set of distinct sifting primes,
$P=\prod_{p\in\mathcal P}p$, and let $D_{\rm nat}\ge2$ be an integer level
such that every $p\in\mathcal P$ is below $D_{\rm nat}$.
The M\"obius function $\mu$ is zero on integers divisible by a prime square
and is $(-1)^k$ on products of $k$ distinct primes, with $\mu(1)=1$;
thus $\mu^2$ is the squarefree indicator on positive integers.
For $d\mid P$ write its prime factors in decreasing order
$d=p_1\cdots p_k$, $p_1>\cdots>p_k$. Define
\[
 \lambda_d^{\pm}=
 \begin{cases}
 \mu(d),&d\mid P,\ d<D_{\rm nat},\ \text{and the corresponding chain tests hold},\\
 0,&\text{otherwise},
 \end{cases}
\]
where the test at position $j$ is
\[
 p_1\cdots p_{j-1}p_j^3<D_{\rm nat}:
 \quad\begin{cases}
 j\text{ even},&\lambda_d^-,\\
 j\text{ odd},&\lambda_d^+.
 \end{cases}
\]
The empty chain passes, so $\lambda_1^\pm=1$.
In real-level consumers the convention $D_{\rm nat}=\lfloor D_{\rm real}\rfloor+1$
makes $d<D_{\rm nat}$ exactly $d\le D_{\rm real}$ for integer $d$;
the chain tests themselves still use $D_{\rm nat}$.
Other natural-level choices, such as the level-six construction below,
are retained as specified by their producers.

The finite Rosser certificates give the pointwise inequalities, for positive
integers $a$,
\[
 \sum_{d\mid\gcd(a,P)}\lambda_d^-
 \le \mathbf1_{\gcd(a,P)=1}
 \le \sum_{d\mid\gcd(a,P)}\lambda_d^+,
 \qquad |\lambda_d^\pm|\le1.
\]
For a finite sequence $\mathcal A$ of positive integers (counting any labels
with their multiplicity), set
$\mathcal A_d=\{a\in\mathcal A:d\mid a\}$ and
$S(\mathcal A,P)=\#\{a\in\mathcal A:\gcd(a,P)=1\}$.
Write $|\mathcal A_d|=Xg(d)+r_d$, where $X\ge0$ is the main mass,
$g$ is the multiplicative density with $g(1)=1$, and $r_d$ is the discrepancy.
Summing the pointwise inequalities and interchanging finite sums gives
\[
 X\sum_{d\mid P}\lambda_d^-g(d)-\sum_{d\mid P}|\lambda_d^-r_d|
 \le S(\mathcal A,P)\le
 X\sum_{d\mid P}\lambda_d^+g(d)+\sum_{d\mid P}|\lambda_d^+r_d|.
\]
The same argument applies to nonnegative weighted sources. This separates
the analytic density comparison from payment of the distribution remainder.
These combinatorial Rosser weights are not the Selberg optimizer used later.
The lower definition and certificate are in
\sourcefile{MathlibNt/SieveTheory/LinearSieve/FiniteWeights.lean}{finite lower Rosser weights};
the odd-position definition is in
\sourcefile{MathlibNt/SieveTheory/LinearSieve/RosserChains.lean}{upper Rosser chains},
and its certificate is in
\sourcefile{MathlibNt/SieveTheory/LinearSieve/UpperRosserDensity.lean}{finite upper Rosser weights}.

\paragraph{The conditioned distribution bill.}
For Chen's conditioning let $P_N=\prod_{2<p\le N^{1/10},\ p\nmid N}p$.
The relevant modulus is $qd$, not $d$; medium primes satisfy
$N^{1/10}<q\le N^{1/3}$. The weighted producer controls the reduced carrier
\inputleannode{bp:foundation-weighted-bv}
\[
 \sum_{\substack{q\text{ in Chen's medium range}\\q\nmid N}}
 \ \sum_{\substack{d\mid P_N\\d<\lfloor N^{1/2-\epsilon}/q\rfloor+1}}
 3^{\omega(d)}\max_{(a,qd)=1}
 \left|\pi(N;qd,a)-\frac{L_*(N)}{\varphi(qd)}\right|
 \ll_{\epsilon,A}\frac{N}{(\log N)^A}.
\]
Here $0<\epsilon<1/6$ and $A>0$.
The proof checks injectivity of $(q,d)\mapsto qd$ on this carrier and pays the
weight by finite Cauchy and divisor-moment estimates, using ordinary BV with
saving exponent $2A+10$.
Nonreduced lanes $q\mid N$ and removed small primes have separate corrections.
This producer enters Chen's conditioned upper sieve; ordinary BV also pays the
base lower-sieve remainder and Li--Liu's strict-endpoint $S_1$ remainder.

\subsection{Constructed sieve functions and uniform comparison}

Write $F$ and $f$ for the constructed dimension-one upper and lower
Jurkat--Richert functions. Their initial values are $F(s)=2e^\gamma/s$
for $0<s\le3$ and $f(s)=0$ for $0<s\le2$; the delay equations are
$(sF(s))'=f(s-1)$ and $(sf(s))'=F(s-1)$ beyond the initial intervals.
For a real level $D_{\rm real}>0$ and cutoff $z>1$, the sieve coordinate is
$s=\log D_{\rm real}/\log z$.

The hats needed for the error estimate are positive extensions of the
normalized derivatives $-F'/(2e^\gamma)$ and $f'/(2e^\gamma)$, rather
than additional counting functions. With $c_\gamma=2e^\gamma$, they are,
for $s>0$,
\[
 \widehat T^+(s)=
 \begin{cases}s^{-2},&s\le3,\\
 \bigl(F(s)-f(s-1)\bigr)/(c_\gamma s),&s>3,
 \end{cases}
 \qquad
 \widehat T^-(s)=
 \begin{cases}2s^{-2},&s\le2,\\
 \bigl(F(s-1)-f(s)\bigr)/(c_\gamma s),&s>2.
 \end{cases}
\]
The delay construction proves positivity, continuity, the weighted delay
relations, and decay, with Suzuki's initial parameter $\widehat\beta=2$:
\inputleannode{bp:foundation-jr-hats}
The definitions and all these properties are supplied together in
\sourcefile{MathlibNt/SieveTheory/LinearSieve/JurkatRichert/JurkatRichert1965Section13HatSource.lean}{constructed hats and their source theorem}.

To state the lower comparison precisely, fix real parameters $d,\Delta,\Theta$
(the letter $d$ here is an error exponent, not a sifting divisor) satisfying
\[
 0<\Delta<1,\qquad d>\frac7{1-\Delta},\qquad \Theta>0,
 \qquad \frac2d<\frac1\Theta,\qquad
 \frac2\Theta+\frac3d<1-\Delta.
\]
They are chosen before the level grows. For $D>1$ define the source cutoff
and the depth-$m$ error envelope by
\[
 \sigma_{\rm src}(D;d)=(\log D)^{1/d}\log\log(27D),
 \qquad
 E_m(D,s;d)=\left(1+\frac{s^d}{\log D}\right)^s s
 \begin{cases}\widehat T^+(s),&m\text{ odd},\\
 \widehat T^-(s),&m\text{ even}.
 \end{cases}
\]
For a bounding sieve $S$, let $\mathcal P_S$ denote its finite sifting prime
set and let $g_S$ denote its multiplicative density, with $0\le g_S(p)<1$.
Put $\mathcal P_S(z)=\{p\in\mathcal P_S:p<z\}$,
$P_S(z)=\prod_{p\in\mathcal P_S(z)}p$, and
$V_S(z)=\prod_{p\in\mathcal P_S(z)}(1-g_S(p))$.
Its dimension-one local-product condition with constant $K$ is
\[
 \prod_{\substack{p\in\mathcal P_S\\z_1\le p<z_2}}
 (1-g_S(p))^{-1}
 \le\frac{\log z_2}{\log z_1}\left(1+\frac K{\log z_1}\right)
 \qquad(2\le z_1\le z_2).
\]
For an integer level $D_{\rm nat}\ge2$ and $4\le s\le6$, use the actual
rounded cutoff and carrier-dependent even depth
\[
 z=\left\lceil D_{\rm nat}^{1/s}\right\rceil,\qquad
 m=2\bigl(|\mathcal P_S(z)|+1\bigr).
\]
The source-coordinate relation is the displayed ceiling; after rounding,
one must not replace $s$ by $\log D_{\rm nat}/\log z$.
Assume $z\ge2$ and $s\le\sigma_{\rm src}(D_{\rm nat};d)$.
With the constructed hats and the fixed parameters above, the lower
comparison chooses its constants before $S$, $K$, $D_{\rm nat}$ and $s$ vary:
\inputleannode{bp:foundation-suzuki-lower}
In particular, there is $C\ge3$ such that for each $K\ge2$ satisfying the
local-product condition,
\[
 V_S(z)\bigl(f(s)-\mathcal E_m(D_{\rm nat},s;d,\Delta,K)\bigr)
 \le \sum_{e\mid P_S(z)}\lambda_e^-g_S(e),
\]
where the weights have level $D_{\rm nat}$ and the \emph{complete} loss is
\[
 \mathcal E_m(D,s;d,\Delta,K)
 =C\exp(\sqrt K)\,E_m(D,s;d)(\log D)^{-\Delta}.
\]
The parameter inequalities, cutoff, envelope and full comparison can be
read respectively in
\sourcefile{MathlibNt/SieveTheory/LinearSieve/Suzuki/SuzukiClaim145SourceParameters.lean}{Suzuki's parameter conditions},
\sourcefile{MathlibNt/SieveTheory/LinearSieve/Suzuki/SuzukiCaseIISourceSigmaDecay.lean}{the source cutoff},
\sourcefile{MathlibNt/SieveTheory/LinearSieve/Suzuki/SuzukiLemma144ErrorObjects.lean}{the exact error envelope}, and
\sourcefile{MathlibNt/SieveTheory/LinearSieve/Suzuki/SuzukiLemma147JurkatRichertInterval.lean}{the uniform lower comparison}.

For Chen the source verifies $(d,\Delta,\Theta)=(16,1/2,7)$ and takes
\[
 D_{\rm nat}=\lfloor N^{1/2-\epsilon}\rfloor+1,
 \qquad s=5-10\epsilon,\qquad 0\le\epsilon<1/10.
\]
For each fixed $\epsilon$ and each $\rho>0$, it then selects an $N$ threshold
for both the source-cutoff condition and $\mathcal E_m<\rho$, for \emph{all}
depths $m$. Indeed, with $s$ fixed the perturbation
$(1+s^{16}/\log D)^s$ tends to $1$ and $(\log D)^{-1/2}$ tends to zero;
depth only chooses between two fixed hat values. A growing prime set and
hence a growing adaptive depth therefore cause no extra loss.
The actual finite-source comparison is
\sourcefile{MathlibNt/SieveTheory/LinearSieve/Suzuki/SuzukiChenVaryingFamilyLowerDensity.lean}{Chen's varying-family lower density},
using \sourcefile{MathlibNt/SieveTheory/LinearSieve/Suzuki/SuzukiChenAdaptiveDepthAbsorption.lean}{all-depth error absorption}
and \sourcefile{MathlibNt/SieveTheory/LinearSieve/Suzuki/SuzukiChenParameterBridge.lean}{the floor/ceiling geometry}.
Li--Liu's beta branch also fixes its coordinate before choosing the integer
threshold; its final fixed-$s$ range is $4\le s<33/8$.

The modern upper comparison gives, for each local-product constant $K>1$
and tolerance $\rho>0$, a cutoff $z_0$ uniform in the bounding sieve:
\inputleannode{bp:foundation-suzuki-upper}
If $z\ge\max(2,z_0)$, every sifting prime is at most $z$,
$D_{\rm real}>0$, and $s=\log D_{\rm real}/\log z$ satisfies
$3/2\le s\le4$, then
\[
 \sum_{e\mid P}\lambda_e^+g_S(e)
 \le (F(s)+\rho)\prod_{p\mid P}(1-g_S(p)),
 \qquad D_{\rm nat}=\lfloor D_{\rm real}\rfloor+1,
\]
where $P=\prod_{p\in\mathcal P_S}p$ and the product on the right runs over
primes. This comparison uses the local-product condition just defined.
The source-cutoff and odd-depth error are absorbed in the choice of $z_0$;
see \sourcefile{MathlibNt/SieveTheory/LinearSieve/Suzuki/SuzukiUpperRosserDensityEndpointDirect.lean}{the upper comparison through four}.
A finite source bound still requires the separate remainder payment derived
from the pointwise Rosser inequality.

\paragraph{Which comparisons the Li--Liu terms actually use.}
The following three applications distinguish the ranges before the explicit
sieve-function formulas are used in the Li--Liu chapter.
\begin{enumerate}
\item Chen's conditioned upper sieve and Li--Liu's $S_2$ and $\mathcal B_{10}$
upper sieves use the preceding $3/2\le s\le4$ theorem. Li--Liu's second
positive $S_1$ term uses the fixed-$s$ beta lower branch; continuity selects
$s<33/8$ before the threshold to approximate its endpoint coefficient.
\item The first positive $S_1$ term instead uses the genuine level-six
construction at $z_N=N^{4/53}$:
\[
 Z_6=\lceil z_N\rceil,\qquad D_6=Z_6^6,\qquad s=6.
\]
For every $\rho>0$, the lower Rosser main sum at this natural level is at
least $(f(6)-\rho)V_N(Z_6)$ eventually, uniformly in the source parameter
$0<\epsilon<1$. Since $D_6\sim N^{24/53}$ and $24/53<1/2$, this choice leaves room below
the ordinary-BV square-root level for the distribution payment.
The definition and density theorem are
\sourcefile{MathlibNt/SieveTheory/LiLiuGoldbachS1LevelSix.lean}{the level-six geometry}
and \sourcefile{MathlibNt/SieveTheory/LiLiuGoldbachS1LowerDensitySix.lean}{the actual level-six lower density}.
Combining this with the main mass, the lower product, and the paid ordinary-BV
remainder gives, for each fixed $0<\epsilon<1$ and $\delta>0$,
\[
 \left(\frac{53}{2}e^{-\gamma}(1-\epsilon)f(6)-\delta\right)\mathcal X_N
 \le S_1(z_N)
\]
eventually for even $N$, where $S_1(z_N)$ denotes the base sift of the
$\epsilon$-truncated difference source at cutoff $z_N$.
The actual assembly is
\sourcefile{MathlibNt/SieveTheory/LiLiuGoldbachS1AlphaNormalizedLower.lean}{the first positive normalized term};
it uses the separately paid level-six error, not the beta endpoint theorem.
The notation $f_*$ in the Li--Liu chapter is this $f$ on $[4,6]$.
\item The $S_3$ upper main sum feeding $G_4,G_5$ uses a \emph{different}
producer. Let $F_S$ denote the constructed continuous Suzuki upper factor;
its subscript means Suzuki, not dependence on the particular sieve $S$.
For each $K>1$ and $\rho>0$, its own uniform cutoff gives
\[
 \sum_{e\mid P}\lambda_e^+g_S(e)
 \le (F_S(s)+\rho)\prod_{p\mid P}(1-g_S(p))
 \qquad(3/2\le s\le6),
\]
with the same real-level, prime-support and local-product conditions as above.
This is \sourcefile{MathlibNt/SieveTheory/LiLiuGoldbachUpperDensitySix.lean}{the upper comparison through six};
its odd-depth error is controlled uniformly in $s$ on this larger interval.
The theorem \sourcefile{MathlibNt/SieveTheory/LiLiuGoldbachS3RosserMainUpper.lean}{the actual S3 Rosser main sum}
supplies the Goldbach local-product constant and specializes the sum to
$\sum_{e\mid P}\lambda_e^+/\varphi(e)$.
It is consumed in
\sourcefile{MathlibNt/SieveTheory/LiLiuGoldbachS3NormalizedUpper.lean}{the normalized S3 upper bound}.
This covers the whole required $53/24\le s\le45/8$ interval, including
its part beyond $4$. It is not a change of the range of the preceding node.
\end{enumerate}

\subsection{Switched distribution is an aggregate estimate}

For a bounded real coefficient $a(m)$ on positive integers, an interval
$I=(A_1,A_2]\cap\mathbb N$ with integer $0\le A_1\le A_2$, an endpoint
$Y\ge0$, a modulus $q\ge1$, and a reduced residue $l$, put
\[
 R_{a,I,\kappa}(Y;q,l)=
 \sum_{\substack{m\in I\\(m,q)=1}}a(m)
 \left(\#\{p:mp\le Y,\ mp\equiv l\pmod q\}
       -\frac{L_\kappa(Y/m)}{\varphi(q)}\right).
\]
The counted variable $p$ is prime. Set
$M_{a,I,\kappa}(Y;q)=\max_{(l,q)=1}|R_{a,I,\kappa}(Y;q,l)|$.
The sum over $m$ stays \emph{inside} the absolute value.
Replacing this by the sum of absolute errors for each $m$ would demand a
stronger, different distribution theorem.

For Liu's particular prime-pair characteristic coefficient and prescribed
source interval, the unweighted Pan specialization is proved at $\kappa=2/\log2$.
\inputleannode{bp:foundation-liupan}
Its arbitrary logarithmic saving is transferred to the
$\mu(q)^2 3^{\omega(q)}$-weighted source by finite weight payment.
The strict source cutoff is only enlarged or embedded with a proved inclusion;
transport to the closed-floor consumer spends a logarithm, replacing $B$ by $B+1$.
A separate additive-normalization estimate gives the canonical $L_2$ coprime
input used by both public Chen implementations.
\inputleannode{bp:foundation-liupan-canonical}

Li--Liu's changing source coefficients require a broader producer, not merely
Chen's fixed characteristic coefficient.
\inputleannode{bp:foundation-pan-bounded}
For every $U>0$ it chooses $C,B,N_0$ before $N$, $I$, and $a$,
assuming $|a(m)|\le1$, $A_2\le N^{2/3}$, and $(\log N)^{2B}\le A_1$,
and bounds $\sum_{q\le Q_B(N)}M_{a,I,*}(N;q)$ by $CN/(\log N)^U$.
The $S_2$ switched remainder is identified with such an aggregate.
For $B_{10}$, the finite window is a difference of prefixes at $N$ and
$\lfloor\epsilon N\rfloor$; both endpoints must satisfy the producer's geometry.
Removing the main-term coprimality gate leaves a signed deleted-main term,
whose absolute contribution is paid separately before the common-mass sieve.

\subsection{Selberg optimization and the Chen triple penalty}

Put $R=\lfloor N^{1/4-\epsilon/2}\rfloor$ and define Liu's paper modulus by $Q=\prod_{p\le R,\ p\nmid N}p$, with the product over primes.
For coefficients supported on $d\mid Q$, $d\le R$, the Selberg square expands
into a main quadratic form plus a signed remainder indexed by pairs of divisors.
The explicit optimizer is normalized at $1$ and has $|\lambda_d|\le1$
on the admissible even source. Its denominator is
\[
 G(N,\epsilon)=\sum_{\substack{d\mid Q\\d\le R}}
                  \prod_{p\mid d}\frac1{p-2}.
\]
For fixed $\epsilon<1/2$, the proved even-integer limit is
\inputleannode{bp:foundation-selberg-denominator}
\[
 \frac{G(N,\epsilon)\mathfrak S_{\mathrm{Liu}}(N)}{\log N}
 \longrightarrow \frac{1/4-\epsilon/2}{2}.
\]
Together with the genuine-logarithmic-integral prime-pair main mass, this gives
the optimized main-term bound used in the Chen upper penalty.

Grouping the remainder by $d=[d_1,d_2]=\operatorname{lcm}(d_1,d_2)$,
the least common multiple, uses $d\le d_1d_2\le R^2$.
For squarefree $d$, each prime has three possible placements in the pair,
which explains the factor $3^{\omega(d)}$.
\inputleannode{bp:foundation-selberg-remainder}
The resulting full-distribution majorant is supplied by the canonical coprime
Pan theorem together with the proved noncoprime correction.
The even Selberg assembly is then transported to corrected Chen triples;
modulus-dividing residual primes and source-boundary fibres are not discarded.
This is Chen's Selberg route, not a substitute for Li--Liu's linear upper sieves.

\subsection{What is shared, and where the arguments separate}

The shared foundation is now visible: PNT and Mertens normalization, ordinary BV,
constructed Jurkat--Richert functions, uniform Suzuki comparisons, and the
low/high-conductor machinery underlying the two Pan producers.
Li--Liu's $S_1$ normalization combines paid ordinary-BV error, lower density,
and the genuine lower sieve product; its $S_2$ and $B_{10}$ normalizations use
bounded-coefficient Pan distribution, upper density, and the upper product.
Their final elementary error absorption is not the source of these inputs.

There is also a distinct well-factorable rectangular distribution lane.
Fix integers $i,j,A\ge0$, a scale constant $C_{\rm scale}\ge1$, and a
real gap $\zeta>0$. These are fixed \emph{before} choosing the threshold
in $x$. For each $x$ above that threshold, the estimate is uniform in
rectangles, residues and coefficients satisfying
\[
 M\ge1,\quad T\ge1,\quad 4MT=x,\quad T=x^\nu,
 \quad \zeta\le\nu\le\frac1{10}+\frac\zeta{10},
 \quad N\in\mathbb N,\quad 0<N\le C_{\rm scale}x.
\]
Choose a finite long-variable set $\mathcal U\subset[M,2M]\cap\mathbb N$
and a short interval $J=(u,v]\cap\mathbb N$ with $T\le u\le v\le2T$.
The long coefficient $\alpha$ satisfies $|\alpha(m)|\le\tau_i(m)$ on
$\mathcal U$, where $\tau_i$ counts ordered factorizations into $i$ positive
integer factors ($\tau_0$ is the indicator of $1$).
The short coefficient is the actual Goldbach prime mask
$\beta_N(n)=\mathbf1_{n\text{ prime},\,(n,N)=1}$.
Set $Q=x^{(5-5\nu)/9-\zeta}$, and let $c$ be a signed well-factorable
coefficient of order $j$ at level $Q$: $|c(q)|\le\tau_j(q)$, and for every
split $Q=Q_1Q_2$ with $Q_1,Q_2\ge1$, it is a Dirichlet convolution
$c=c_1*c_2$ of coefficients supported on $1\le n\le Q_1$ and
$1\le n\le Q_2$, respectively, each bounded in absolute value by $\tau_j(n)$.
The centered signed error is
\[
 \mathcal E_{\rm rect}=
 \sum_{\substack{1\le q\le\lfloor Q\rfloor\\(q,N)=1}}c(q)
 \left(
 \sum_{\substack{m\in\mathcal U,\ n\in J\\mn\equiv N\pmod q}}
       \alpha(m)\beta_N(n)
 -\frac1{\varphi(q)}
 \sum_{\substack{m\in\mathcal U,\ n\in J\\(mn,q)=1}}
       \alpha(m)\beta_N(n)
 \right).
\]
\inputleannode{bp:foundation-fouvry-rectangle}
The theorem gives $|\mathcal E_{\rm rect}|\le x/(\log x)^A$ for all these
choices after the one threshold depending only on $i,j,A,C_{\rm scale},\zeta$.
In particular the Goldbach residue $N$ may vary in the displayed range,
whereas the positive gap $\zeta$ is not allowed to shrink with $x$.
The precise interval geometry is in
\sourcefile{MathlibNt/AnalyticNumberTheory/LargeSieve/LiLiuFouvryPrimeC2.lean}{the prime rectangle},
the convolution condition in
\sourcefile{MathlibNt/AnalyticNumberTheory/LargeSieve/LiLiuPrereqFouvryFactorConvolution.lean}{signed well-factorability},
and the centered error in
\sourcefile{MathlibNt/AnalyticNumberTheory/LargeSieve/LiLiuPrereqFouvryDispersion.lean}{the bilinear discrepancy}.
The parameter order and the literal residue $N$ are those of
\sourcefile{MathlibNt/AnalyticNumberTheory/LargeSieve/LiLiuFouvryKGoldbachRectangle.lean}{the Goldbach rectangle theorem}.
Li--Liu's flexible $G_{12}$ rectangles consume this producer unchanged;
curved-region boundaries and exceptions are paid separately in that chapter.

\subsection{Inside the distribution proofs}
A Dirichlet character modulo $q$ is a multiplicative function on the reduced
residue classes, extended by zero to nonunits. Character orthogonality
converts residue-class errors into character sums. The conductor is the
smallest modulus from which a character is induced; passing to primitive
characters also introduces cofactor sums and their totient weights.
The principal character is the character equal to one on all units.

For ordinary Bombieri--Vinogradov, three branches meet before the maximal
prime-counting estimate is obtained.
\begin{enumerate}
\item The principal part is paid by the quantitative prime number theorem,
with the actual logarithmic-integral normalization and prime-power correction.
\item Nonprincipal characters of small conductor use Siegel--Walfisz
estimates: arbitrary inverse-logarithmic saving while the conductor stays
in the prescribed logarithmic range. The implementation constructs the
smoothing function and obtains the required uniform quadratic $L$-value
lower bound from its proved Landau--Siegel theorem.
\inputleannode{bp:foundation-bv-low}
\item Large conductors use Vaughan's identity to split the von Mangoldt
weight into Type I, Type II and a small remainder. Type I retains a long
summation variable and uses periodic character cancellation with Cauchy;
Type II is bilinear and uses the all-aspect large-sieve estimate on the
chosen high-conductor family. One common choice of the two Vaughan
cutoffs is used in all three estimates. Prefix maxima and conductor
sums are paid before the resulting inverse-logarithmic bound is consumed.
\inputleannode{bp:foundation-bv-high}
\end{enumerate}
The unconditional BV producer combines these branches and passes to prime
counts by partial summation. Its proof is a modern implementation of the
needed distribution estimate, with each branch's actual inputs supplied.

For Pan distribution, the object throughout the proof is the signed
coefficient--prime convolution defined above. The low branch uses the
nonprincipal prime-prefix Siegel--Walfisz estimate. The high branch keeps
the short and long coefficient polynomials together; a Perron integral
separates the hyperbolic product constraint while preserving their
character coupling. Low and high estimates hold with the cofactor still
free after the common threshold. The primitive-conductor assembly then
pays the reciprocal-totient cofactor sum, and a separate principal estimate
completes the aggregate. This explains why ordinary prime BV cannot
simply replace the Pan input in a switched sieve.

The well-factorable branch keeps cancellation also in the modulus
coefficient: a level-$Q$ well-factorable coefficient admits appropriately
supported convolution factorizations when $Q$ is split into two levels.
The signed version and its divisor bounds are part of the implemented
rectangle theorem's assumptions. Applying its bilinear estimates to each
admissible rectangle is followed by the separate boundary and exception
payments in the Li--Liu chapter.
The subsequent $G_{67},G_9,G_{11},G_{12}$ integral and signed-count assemblies
belong to the Li--Liu argument, rather than being silently absorbed into a
claim that the common analytic estimates alone prove both endpoints.
